\section{引言：从单 GPU 到数据中心}

上一节课（Lecture 5/6）我们深入探讨了单 GPU 的工作原理和性能优化。然而，当代大语言模型的规模已经远远超出了单 GPU 的能力边界——无论是计算量还是内存容量。本节课的核心问题是：\textbf{如何将训练工作高效地分配到成百上千个 GPU 上？}

\begin{importantbox}{本课三大核心目标}
\begin{enumerate}
    \item 理解多机并行训练的必要性——计算和内存两方面的驱动力
    \item 掌握四种主要并行策略——数据并行、流水线并行、张量并行、序列并行
    \item 学会组合使用——3D/4D 并行的设计经验法则
\end{enumerate}
\end{importantbox}

\begin{figure}[H]
\centering
\includegraphics[width=0.95\textwidth]{slides-images/slide-001.jpg}
\caption{课程大纲：从网络基础到并行策略再到实际案例\protect\footnotemark}
\end{figure}
\footnotetext{Slides 第2页。}

\subsection{计算驱动：超级计算机的规模}

单 GPU 的算力虽然在飞速增长（超指数增长曲线），但如果我们希望\textbf{此时此刻}训练最强大的语言模型，就必须依赖多机并行。当今世界最快的超级计算机拥有 exaFLOPS 级别的总算力，这正是训练最大模型所需要的资源规模。

\begin{figure}[H]
\centering
\includegraphics[width=0.95\textwidth]{slides-images/slide-002.jpg}
\caption{GPU 算力增长与超级计算机规模对比\protect\footnotemark}
\end{figure}
\footnotetext{Slides 第3页。}

\subsection{内存驱动：模型已超出单卡容量}

除了计算之外，内存也是核心瓶颈。大模型拥有数十亿甚至数千亿参数，加上优化器状态、梯度和激活值，所需内存远超单 GPU 的 HBM 容量。

\begin{figure}[H]
\centering
\includegraphics[width=0.95\textwidth]{slides-images/slide-003.jpg}
\caption{模型参数量增长与 GPU 内存增长的对比\protect\footnotemark}
\end{figure}
\footnotetext{Slides 第4页。}

\begin{knowledgebox}{新的计算单元：数据中心}
从这节课开始，我们的基本计算单元不再是单个 GPU，而是整个\textbf{数据中心}。我们需要设计算法和分片策略来实现两个目标：
\begin{itemize}
    \item \textbf{线性内存扩展}：GPU 数量翻倍，可训练的最大模型规模也翻倍
    \item \textbf{线性计算扩展}：GPU 数量翻倍，有效计算吞吐量也翻倍
\end{itemize}
\end{knowledgebox}

\subsection{本章小结}

多 GPU 并行训练的必要性来自两个方面：（1）计算需求——单 GPU 算力不足以在合理时间内完成训练；（2）内存需求——模型参数、优化器状态等无法放入单卡。本课将介绍如何通过不同的并行策略，在数据中心规模上同时解决这两个问题。

%% ========== Part 2: 网络层次与集合通信 ==========

